\section{Графы и их представление в ЭВМ}
\label{sec:graphs-representation}

\subsection{Основные определения}

Графом называется пара
\[
 G=(V,E),
\]
где $V$ --- множество вершин, а $E$ --- множество рёбер. В простом
неориентированном графе ребро --- неупорядоченная пара $\{u,v\}$ различных
вершин. В ориентированном графе дуга --- упорядоченная пара $(u,v)$.

Допускаются также мультиграфы (кратные рёбра), псевдографы (петли), взвешенные
графы и помеченные графы. Далее, если не оговорено иное, речь идёт о конечном
простом графе. Обозначим
\[
 n=|V|,\qquad m=|E|.
\]

Вершины $u,v$ смежны, если соединены ребром. Ребро инцидентно каждой из своих
концевых вершин. Степень вершины
\[
 d(v)=|\{e\in E: v\in e\}|.
\]
Для орграфа различают
\[
 d_{\rm out}(v),\qquad d_{\rm in}(v).
\]

\textbf{Лемма о рукопожатиях.}
Для неориентированного графа
\[
 \sum_{v\in V}d(v)=2m.
\]
Действительно, каждое ребро учитывается в сумме степеней ровно дважды. Следствие:
число вершин нечётной степени всегда чётно.
Для орграфа
\[
 \sum_v d_{\rm out}(v)=\sum_v d_{\rm in}(v)=m.
\]

\subsection{Подграфы, маршруты и связность}

Подграф $H=(V_H,E_H)$ графа $G$ удовлетворяет $V_H\subseteq V$ и
$E_H\subseteq E$. Если $V_H=V$, подграф называется остовным. Подграф,
порождённый множеством вершин $W\subseteq V$, содержит все рёбра исходного графа
с обоими концами в $W$.

Маршрут --- последовательность смежных рёбер/вершин. Если рёбра не повторяются,
получаем цепь (trail), если не повторяются вершины --- простой путь. Замкнутая
простая цепь называется циклом.

Две вершины связаны, если существует соединяющий их путь. Это отношение
эквивалентности; его классы --- компоненты связности. Для связного графа расстояние
$d(u,v)$ определяется как минимальное число рёбер в пути из $u$ в $v$.

Эксцентриситет, радиус и диаметр:
\[
 e(v)=\max_{u\in V}d(u,v),\qquad
 R(G)=\min_v e(v),\qquad
 D(G)=\max_v e(v).
\]

\subsection{Матрица смежности}

Для нумерованных вершин $v_1,\ldots,v_n$ матрица смежности
$A=(a_{ij})$ задаётся формулой
\[
 a_{ij}=\begin{cases}
 1,& (v_i,v_j)\in E,\\
 0,&\text{иначе.}
 \end{cases}
\]
Для простого неориентированного графа $A=A^T$ и $a_{ii}=0$.
Для взвешенного графа вместо $1$ часто хранят вес ребра, хотя при этом нужно
отдельно выбрать значение, обозначающее отсутствие ребра.

Главное достоинство --- проверка наличия ребра за $O(1)$. Недостаток --- память
$O(n^2)$ независимо от числа рёбер; перечисление всех соседей вершины занимает
$O(n)$.

Матрица смежности имеет важный комбинаторный смысл:
\[
 (A^k)_{ij}
\]
равно числу маршрутов длины $k$ из $v_i$ в $v_j$. Это доказывается индукцией по
$k$: умножение матриц суммирует все способы выбрать предпоследнюю вершину.

Для булевой арифметики из матриц $I,A,A^2,\ldots$ можно получать отношение
достижимости. На практике транзитивное замыкание матрицы смежности можно строить,
например, алгоритмом Уоршелла за $O(n^3)$.

\subsection{Матрица инцидентности}

Матрица инцидентности имеет размер $n\times m$. Для простого неориентированного
графа в столбце каждого ребра стоят две единицы. Для ориентированного графа удобно
использовать знаковую матрицу
\[
 b_{ve}=\begin{cases}
 -1,&e\text{ выходит из }v,\\
 +1,&e\text{ входит в }v,\\
 0,&\text{иначе.}
 \end{cases}
\]
Каждый столбец тогда имеет сумму ноль.

Плотное хранение требует $O(nm)$ памяти, поэтому в вычислительных задачах такая
матрица обычно хранится в разреженном формате. Инцидентностные матрицы естественно
возникают в сетевых задачах, дискретных законах сохранения и при построении
графового лапласиана.

Для неориентированного графа ориентируем каждое ребро произвольно и построим
знаковую матрицу $B$. Тогда
\[
 L=BB^T=D-A
\]
--- матрица Лапласа графа, где $D=\operatorname{diag}(d(v_1),\ldots,d(v_n))$.
В частности,
\[
 x^TLx=\sum_{\{i,j\}\in E}(x_i-x_j)^2\ge0,
\]
то есть $L$ положительно полуопределена. Этот факт часто используется в
спектральных алгоритмах на графах.

\subsection{Списки смежности}

Для каждой вершины хранится список её соседей. В неориентированном графе каждое
ребро появляется в двух списках, поэтому общий объём памяти
\[
 O(n+m).
\]
Перебор соседей вершины занимает $O(d(v))$, а полный просмотр всех списков ---
$O(n+m)$. Поэтому списки смежности --- стандартное представление разреженных
графов и естественная структура для BFS и DFS.

Проверка конкретного ребра $(u,v)$ в неупорядоченном списке требует
$O(d(u))$. Возможны улучшения:
\begin{itemize}
 \item хранить соседей в хеш-таблице --- ожидаемая проверка $O(1)$;
 \item хранить отсортированно --- проверка $O(\log d(u))$ бинарным поиском;
 \item использовать сбалансированное дерево --- $O(\log d(u))$ и удобные
       динамические обновления.
\end{itemize}

\subsection{CSR/CSC --- компактные разреженные форматы}

Для больших статических графов списки смежности часто реализуют без отдельных
объектов-списков. В формате CSR (Compressed Sparse Row) используются два массива:
\begin{itemize}
 \item $adj$ --- подряд записанные все соседи;
 \item $offset[v]$ --- начало диапазона соседей вершины $v$.
\end{itemize}
Соседи $v$ находятся в
\[
 adj[offset[v] \ldots offset[v+1]).
\]
Память остаётся $O(n+m)$, данные расположены компактно и хорошо используют
кэш-память процессора. Для ориентированного графа аналогичный CSC-формат удобен
для хранения входящих дуг.

\subsection{Список рёбер}

Граф хранится массивом
\[
 (u_e,v_e)
\quad\text{или}\quad
 (u_e,v_e,w_e),\qquad e=1,\ldots,m.
\]
Память $O(m)$. Это представление удобно в алгоритмах, которые сортируют или
последовательно обрабатывают рёбра: Краскал, Беллман--Форд, построение потоковой
сети. Но перечисление соседей отдельной вершины без индекса требует $O(m)$.

\subsection{Сравнение представлений}

\begin{center}
\renewcommand{\arraystretch}{1.2}
\begin{tabular}{|l|c|c|c|}
\hline
Представление & Память & Проверка ребра & Перебор соседей $v$\\ \hline
Матрица смежности & $O(n^2)$ & $O(1)$ & $O(n)$\\ \hline
Список смежности & $O(n+m)$ & $O(d(v))$ & $O(d(v))$\\ \hline
CSR & $O(n+m)$ & обычно $O(d(v))$ & $O(d(v))$\\ \hline
Список рёбер & $O(m)$ & $O(m)$ & $O(m)$\\ \hline
Матрица инцидентности & $O(nm)$ плотн. & --- & ---\\ \hline
\end{tabular}
\end{center}

Отсюда видно, что выбор структуры данных является частью алгоритма. Один и тот же
теоретический алгоритм может иметь существенно разную реальную стоимость при
разных представлениях графа.

\subsection{Что важно сказать на экзамене}

Нужно определить граф, степень, путь и связность, сформулировать лемму о
рукопожатиях, а затем сравнить способы хранения. Для плотного графа естественна
матрица смежности, для разреженного --- списки смежности/CSR. Важно знать оценки
памяти и времени основных операций и комбинаторный смысл степеней матрицы
смежности.

\newpage
